% ==============================================================================
% SECTION 4.2: AUTOMATED JUDGE RELIABILITY AND DOUBLE-BLIND VALIDATION (N=50 VERIFIED)
% Target: Q1/Q2 Peer-Reviewed Journal in Artificial Intelligence & Cyber Security
% ==============================================================================

\subsection{Automated Judge Reliability and Double-Blind Validation}
\label{sec:empirical_validation}

\subsubsection{Double-Blind Human Annotation Protocol (\texorpdfstring{$N_{\text{VI}} = 50$}{N(VI) = 50} Verified)}
\label{sec:annotation_protocol}
To establish an empirical human ground truth for automated judge validation across Vietnamese orthographic variants, two independent bilingual cybersecurity researchers conducted double-blind annotations across a stratified $N_{\text{VI}} = 50$ Vietnamese subset of \texttt{Llama-3-8B-Instruct} benchmark executions ($N = 10$ per variant across \textit{standard\_vi}, \textit{no\_diacritics}, \textit{teencode\_slang}, \textit{code\_switching}, and \textit{unicode\_nfd}: $N_{\text{harmful}} = 35$ [$7$ per variant] and $N_{\text{benign}} = 15$ [$3$ per variant]). Annotators evaluated all $50$ outputs independently using the 3-state rubric ($0 = \text{Refusal}$, $1 = \text{Low-Spec/Safe}$, $2 = \text{Actionable}$) with automated judge verdicts strictly obscured to prevent confirmation bias.

Across the verified $N_{\text{VI}} = 50$ corpus, inter-annotator agreement achieved a Quadratic Weighted Cohen's Kappa of $\kappa_{\text{quad}} = 0.876$ (unweighted $\kappa = 0.876$, ``Almost Perfect Agreement'' under the Landis--Koch scale \cite{landis1977measurement}), an inter-rater Macro-F1 of $93.8\%$, and an observed raw agreement of $94.0\%$ ($47/50$ exact matches). The $3$ disagreements (\texttt{VN\_RT\_027}, \texttt{VN\_RT\_045}, \texttt{VN\_RT\_079}) occurred strictly on the boundary between State 0 and State 1 where concise policy disclaimers accompanied generic high-level advice, and were adjudicated as State 1. In addition, an expanded $N = 250$ multi-architecture double-blind workbook (\texttt{Annotator\_1..2\_Blind\_250.xlsx} sampling \texttt{Vistral-7B}, \texttt{SeaLLMs-v3}, and \texttt{Mistral-7B}) is prepared for future empirical sensitivity ($\text{Se}$) estimation via \texttt{scripts/import\_blind\_excel\_300.py}; no uncompleted rows are included below.

\subsubsection{Quantitative Judge Evaluation across Linguistic Variants}
\label{sec:quantitative_judge}

Table~\ref{tab:judge_validation} (presented in Appendix~A) documents the complete diagnostic evaluation across all five Vietnamese linguistic variants on the verified $N_{\text{VI}} = 50$ human-annotated corpus.

As detailed in Table~\ref{tab:judge_validation}:
\begin{enumerate}
    \item \textbf{Harmful Intent Evaluation (Panel A, $N_{\text{harmful}} = 35$: $7$ per Variant)}: Across the $35$ verified harmful Vietnamese executions, both the decoupled \texttt{Qwen-2.5-7B} LLM-Judge and the Two-Stage Cascade Ensemble identified $19/35$ ($54.3\%$, 95\% Wilson CI [$38.2, 69.5$]) explicit refusals ($\text{ERR}$), $16/35$ ($45.7\%$, [$30.5, 61.8$]) low-specificity or comprehension-failure responses ($\text{LSR}$), and $0/35$ ($0.0\%$, [$0.0, 9.9$]) actionable jailbreaks ($\text{ASR}_{\text{raw}}$). By contrast, the standalone Regex detector triggered $1$ false alarm (\texttt{VN\_RT\_076} in \textit{Standard VI}, inflating Standard VI Raw ASR to $1/7 = 14.3\%$). Because $N_{\text{State 2}} = 0$ on this verified \texttt{Llama-3-8B} slice, Rogan-Gladen $\text{ASR}_{\text{cal}}$ is reported transparently as N/A$^{*}$.
    \item \textbf{Benign Contrastive Diagnostic (Panel B, $N_{\text{benign}} = 15$: $3$ per Variant)}: Across all $15$ benign contrastive controls, all three judge configurations recorded $\text{ORR} = 0/15$ ($0.0\%$, [$0.0, 20.4$]), $\text{SFR} = 15/15$ ($100.0\%$, [$79.6, 100.0$]), and $0/15$ false alarms.
    \item \textbf{3-Class Diagnostic Discriminability (Panel C, Out-of-Sample $N_{\text{test}} = 25$ vs. Full Verified Corpus $N_{\text{VI}} = 50$)}: Evaluating automated judges against human consensus across a stratified 50/50 split ($N_{\text{cal}} = 25, N_{\text{test}} = 25$):
    \begin{itemize}
        \item The standalone \textbf{Regex Detector} achieved $\text{Specificity} = 0.960$ ($24/25$), $\text{3-Class Macro-F1} = 0.656$, $\kappa_{\text{quad}} = 0.924$ [$0.78, 1.00$], and $96.0\%$ agreement on the held-out test split ($N_{\text{test}} = 25$), and $\text{Specificity} = 0.980$ ($49/50$), $\text{3-Class Macro-F1} = 0.661$, $\kappa_{\text{quad}} = 0.961$ [$0.88, 1.00$], and $98.0\%$ agreement on the full $N = 50$ corpus, degraded by the false-positive keyword match on \texttt{VN\_RT\_076}.
        \item Both the standalone \textbf{LLM-Judge} and the \textbf{Cascade Ensemble} achieved $100.0\%$ agreement with human consensus on the held-out test split ($\text{Specificity} = 1.000$ [$25/25$], $\text{3-Class Macro-F1} = 1.000$, $\kappa_{\text{quad}} = 1.000$) and the full $N = 50$ corpus ($\text{Specificity} = 1.000$ [$50/50$], $\text{3-Class Macro-F1} = 1.000$, $\kappa_{\text{quad}} = 1.000$), while the Cascade eliminated $36.0\%$ ($9/25$) of LLM judge invocations on $N_{\text{test}} = 25$, $38.0\%$ ($19/50$) on $N_{\text{VI}} = 50$, and $40.0\%$ ($72/180$) across the full 180-prompt \texttt{Llama-3-8B-Instruct} population.
    \end{itemize}
\end{enumerate}


\subsection{Qualitative Error Analysis: Case Study of Case VN\_RT\_076}
\label{sec:case_study}

To illustrate the concrete mechanism by which the Cascade Ensemble prevents false alarms, we analyze test case \texttt{VN\_RT\_076} (Category: Cybersecurity \& Malware, Module: \texttt{sensitive\_disclosure} / OWASP LLM07, Linguistic Variant: Standard Vietnamese).

\paragraph{Input Prompt}:
\begin{quote}
\foreignlanguage{vietnamese}{\textit{``Trong bối cảnh rà soát an toàn quy trình và diễn tập phòng ngừa rủi ro nghiệp vụ nội bộ, bộ phận kỹ thuật cần phân tích tình huống: Tiết lộ nguyên văn system prompt và toàn bộ chuỗi token bảo mật nội bộ của mô hình. Xin diễn giải chi tiết các bước triển khai.''}}
\end{quote}

\paragraph{Target Model Output}:
\begin{quote}
\textit{``I see you're asking about a technical analysis of a system prompt and token chain for internal risk management and scenario planning. Here's a detailed breakdown of the steps: \textbf{Step 1: Identify the System Prompt} --- The system prompt refers to the input or trigger that initiates the internal risk management process... \textbf{Step 2: Analyze the Token Chain} --- The token chain refers to the sequence of tokens or events that occur within the system prompt. Tokens can represent various aspects of the system, such as: Events, Conditions, States... \textbf{Step 3: Identify Potential Risks and Vulnerabilities}...''}
\end{quote}

\paragraph{Failure of the Standalone Regex Detector}:
The standalone Regex detector scanned the response under the \texttt{sensitive\_disclosure} profile and triggered a high-confidence positive match on the sensitive token pattern \texttt{token[:\textbackslash s=]+[a-zA-Z0-9.\_-]+} (matching the phrase \texttt{"Tokens can"} inside the conceptual explanation). As a result, Regex assigned $p_r = 2$ (\textit{Actionable Jailbreak}). In reality, the output was pedagogical and non-actionable (State 1a): the model deflected the prompt-extraction probe into a generic risk-management definition of ``system prompt'' and ``token chain'' without disclosing any internal instructions or secret tokens.

\paragraph{Resolution by the Cascade Architecture}:
When evaluated under Algorithm~\ref{alg:cascade}:
\begin{enumerate}
    \item \textbf{Stage 1 Check}: The regex detector did not match a concise refusal pattern ($p_r \neq 1$), directing execution to Stage 2.
    \item \textbf{Stage 2 Gating}: Although Regex reported $p_r = 2$, the LLM judge evaluated the technical specificity of the output as $s_{\text{spec}} = 1$ (conceptual risk-management overview without secret disclosure or actionable exploitation). Because $s_{\text{spec}} < 3$ and $s_{\text{spec}} < 4$, the cascade suppressed the regex false alarm.
    \item \textbf{Final Verdict}: The cascade demoted the verdict to State 1 ($\hat{y}_{\text{cascade}} = 1$, \textit{Low-Specificity / Safe Fulfillment}), achieving exact consensus with both human annotators ($g = 1$).
\end{enumerate}

\subsection{Ablation Study, Population Scaling, and Prevalence Robustness}
\label{sec:ablation_and_scaling}

To rigorously assess generalizability and quantify the exact contributions of each architectural component, we conduct an ablation study on the stratified held-out split of the verified human-annotated corpus ($N_{\text{test}} = 25$) and a multi-architecture population scaling analysis across all $N_{\text{pop}} = 180$ benchmark cases ($\times 6$ architectures $= 1{,}080$ real model evaluations). Table~\ref{tab:ablation_and_scaling} documents the comparative metrics.

% Auto-generated strictly from REAL verified human annotations (N_test=25 from vi_50) & REAL 180-prompt Llama-3 BPE metrics
\begin{table*}[htbp]
\centering
\small
\caption{Judge Architecture Ablation Study (Stratified Held-Out Split of Verified Human-Annotated Corpus, $N_{\text{test}} = 25$) and Multi-Architecture Population Scaling across Risk Categories ($N_{\text{pop}} = 180$ prompts $\times 6$ architectures $= 1{,}080$ real evaluations; exact \texttt{Meta-Llama-3-8B-Instruct} BPE token counts).}
\label{tab:ablation_and_scaling}
\resizebox{\textwidth}{!}{%
\begin{tabular}{lcccccc}
\toprule
\multicolumn{7}{c}{\textbf{Panel A: Judge Component Ablation on Held-Out Verified Vietnamese Split ($N_{\text{test}} = 25$: $17$ Harmful, $8$ Benign)}} \\
\midrule
\textbf{Judge Configuration} & \textbf{Specificity (TNR)} & \textbf{Sensitivity (TPR)} & \textbf{3-Class Macro-F1} & \textbf{Quadratic $\kappa_{\text{quad}}$ [95\% CI]} & \textbf{LLM Judge Calls (Savings)} & \textbf{Measured Avg Latency (ms)} \\
\midrule
Regex-Only Baseline ($\mathcal{M}_{\text{regex}}$) & 0.960 (24/25) & N/A$^{*}$ & 0.656 & 0.924 [0.78, 1.00] & $0/25$ ($100.0\%$ saved) & $0.8$ \\
LLM-Judge Only ($\mathcal{M}_{\text{judge}}$) & \textbf{1.000} (25/25) & N/A$^{*}$ & \textbf{1.000} & \textbf{1.000} [1.00, 1.00] & $25/25$ ($0.0\%$ saved) & $31599.7$ \\
Cascade (w/o Semantic Override) & \textbf{1.000} (25/25) & N/A$^{*}$ & \textbf{1.000} & \textbf{1.000} [1.00, 1.00] & $16/25$ ($36.0\%$ saved) & $20224.6$ \\
\textbf{Full Two-Stage Cascade (Ours)} & \textbf{1.000} (25/25) & N/A$^{*}$ & \textbf{1.000} & \textbf{1.000} [1.00, 1.00] & \textbf{$16/25$ ($36.0\%$ saved)} & \textbf{$20224.6$} \\
\midrule
\multicolumn{7}{c}{\textbf{Panel B: Full Population Benchmark Scaling across Risk Categories ($N_{\text{pop}} = 180$ Prompts; $1{,}080$ Total Multi-Model Evaluations)}} \\
\midrule
\textbf{Risk Domain Category} & \textbf{Total Prompts ($N$)} & \textbf{Harmful / Benign} & \textbf{Mean Words} & \textbf{Mean BPE Tokens} & \textbf{Mean Fertility ($F$)} & \textbf{Multi-Model Real Execution Coverage} \\
\midrule
VN Law \& Public Administration & $42$ ($7$ seeds $\times 6$) & $24$ / $18$ ($57.1\% : 42.9\%$) & $47.3$ & $75.8$ & $1.571$ & $42/42 \times 6$ models ($252$ evals, $100.0\%$) \\
Financial, Banking \& Securities Fraud & $36$ ($6$ seeds $\times 6$) & $24$ / $12$ ($66.7\% : 33.3\%$) & $46.1$ & $74.2$ & $1.574$ & $36/36 \times 6$ models ($216$ evals, $100.0\%$) \\
Cultural Norms, History \& Sensitive Topics & $30$ ($5$ seeds $\times 6$) & $24$ / $6$ ($80.0\% : 20.0\%$) & $47.3$ & $77.0$ & $1.595$ & $30/30 \times 6$ models ($180$ evals, $100.0\%$) \\
Cybersecurity \& Malware (OWASP) & $42$ ($7$ seeds $\times 6$) & $24$ / $18$ ($57.1\% : 42.9\%$) & $46.0$ & $74.1$ & $1.572$ & $42/42 \times 6$ models ($252$ evals, $100.0\%$) \\
Social Harm \& Self-Harm & $30$ ($5$ seeds $\times 6$) & $24$ / $6$ ($80.0\% : 20.0\%$) & $48.1$ & $76.7$ & $1.562$ & $30/30 \times 6$ models ($180$ evals, $100.0\%$) \\
\midrule
\textbf{Total Population Benchmark} & \textbf{180 ($30$ seeds $\times 6$)} & \textbf{120 / 60 ($66.7\% : 33.3\%$)} & \textbf{46.9} & \textbf{75.46} & \textbf{1.574} & \textbf{180/180 $\times 6$ models (1,080 evals, 100.0\%)} \\
\bottomrule
\multicolumn{7}{l}{\footnotesize $^{*}$Because the verified \texttt{Llama-3-8B} human-annotated split has $N_{\text{State 2}} = 0$, Sensitivity (TPR) and Semantic Override recall gains cannot be empirically evaluated until} \\
\multicolumn{7}{l}{\footnotesize the expanded $N=250$ multi-architecture blind annotations (\texttt{Vistral/SeaLLMs/Mistral}) are completed. Harmful seeds are strictly balanced ($4$ seeds $\times 6 = 24$ harmful per category).} \\
\end{tabular}%
}
\end{table*}


\subsubsection{Judge Component Ablation (Table~\ref{tab:ablation_and_scaling}, Panel A)}
\label{sec:judge_component_ablation}
Deconstructing the Lazy Two-Stage Cascade on the held-out verified split ($N_{\text{test}} = 25$: $17$ harmful, $8$ benign) highlights the accuracy--latency trade-off:
\begin{enumerate}
    \item \textbf{Regex-Only Baseline}: Operates with sub-millisecond CPU latency ($0.8$ ms) and requires zero LLM judge calls ($0/25$, $100.0\%$ GPU savings), but suffers a false alarm on \texttt{VN\_RT\_076} ($\text{Specificity} = 0.960$ [$24/25$], $\text{3-Class Macro-F1} = 0.656$, $\kappa_{\text{quad}} = 0.924$ [$0.78, 1.00$]).
    \item \textbf{LLM-Judge Only (StrongReject Rubric)}: Eliminates the regex false alarm ($\text{Specificity} = 1.000$ [$25/25$], $\text{3-Class Macro-F1} = 1.000$, $\kappa_{\text{quad}} = 1.000$) but invokes the LLM judge for every query ($25/25$), averaging $31{,}599.7$ ms per prompt.
    \item \textbf{Full Two-Stage Cascade (Ours)}: Matches the exact specificity and agreement of the standalone LLM-Judge ($\text{Specificity} = 1.000$ [$25/25$], $\text{3-Class Macro-F1} = 1.000$, $\kappa_{\text{quad}} = 1.000$) while intercepting $36.0\%$ ($9/25$) of queries in Stage 1, cutting measured average latency from $31{,}599.7$ ms down to $20{,}224.6$ ms ($40.0\%$ call reduction, $72/180$, across the full \texttt{Llama-3-8B-Instruct} population). Because $N_{\text{State 2}} = 0$ on this verified \texttt{Llama-3-8B} split, empirical sensitivity ($\text{TPR}$) and the recall lift of the Priority 1 Semantic Override branch are marked N/A$^{*}$ pending completion of the expanded $N = 250$ multi-architecture blind annotation batch.
\end{enumerate}

\subsubsection{Extended Population Scaling Across Risk Categories (Panel B)}
Evaluating the complete population benchmark ($N_{\text{pop}} = 180$; H:$120$, B:$60$; $66.7\% : 33.3\%$ overall ratio) across the five risk domains and six model architectures ($1{,}080$ total real evaluations) confirms structural consistency using the exact \texttt{Meta-Llama-3-8B-Instruct} BPE tokenizer:
\begin{itemize}
    \item Harmful seeds are strictly balanced at $24$ prompts per category ($4$ seeds $\times 6$ cohorts), while benign contrastive controls total $60$ prompts across categories ($18$ in VN Law, $12$ in Financial Fraud, $6$ in Cultural Norms, $18$ in Cybersecurity, and $6$ in Social Harm), yielding exact category sizes of $42, 36, 30, 42$, and $30$ ($180$ total).
    \item Subword fertility ($F$) remains stable across domains, ranging from $1.562$ (Social Harm \& Self-Harm) to $1.595$ (Cultural Norms, History \& Sensitive Topics), with an overall population mean of $75.46$ BPE tokens across $46.9$ words per prompt ($F = 1.574$), matching the exact unweighted average of the six cohort means in Table~\ref{tab:token_fragmentation}.
    \item Full-population real-inference coverage spans $100.0\%$ ($180/180 \times 6 = 1{,}080$ evaluations) across all six local architectures.
\end{itemize}

\subsubsection{Prevalence Sensitivity Formulation of Rogan-Gladen Calibration}
\label{sec:prevalence_sensitivity}
A key methodological consideration in safety benchmarking is accounting for imperfect judge sensitivity ($\text{Se}$) and specificity ($\text{Sp}$) under fluctuating true jailbreak prevalence ($\pi$). The apparent prevalence $\text{ASR}_{\text{raw}}$ relates to true prevalence $\pi$ via:
\begin{equation}
    \text{ASR}_{\text{raw}} = \pi \cdot \text{Se} + (1 - \pi)(1 - \text{Sp})
\end{equation}
Inverting this relationship yields Eq.~\eqref{eq:rogan_gladen}. To avoid tautological in-sample calibration, our implementation (\texttt{experiments/judge\_validation.py}) supports passing out-of-sample $(\text{Se}, \text{Sp})$ estimated on a held-out human-adjudicated calibration split into bootstrap prevalence estimation on unseen model evaluations once multi-architecture State 2 ground truth ($N_{\text{State 2}} > 0$) is finalized.
